\documentclass[12pt]{article}

\usepackage[utf8]{inputenc}
\usepackage{amsmath, amssymb, amsthm}
\usepackage{geometry}
\geometry{a4paper, margin=2.5cm}

\title{Pi-Sistem ($\pi$-Sistem)}
\author{}
\date{}

\theoremstyle{definition}
\newtheorem{definition}{Definiție}

\begin{document}

\maketitle

\begin{definition}
O familie de mulțimi $\mathcal{P}$ se numește \emph{pi-sistem} (sau $\pi$-sistem) dacă:
\begin{enumerate}
    \item $\Omega \in \mathcal{P}$.
    \item Dacă $A, B \in \mathcal{P}$, atunci $A \cap B \in \mathcal{P}$.
\end{enumerate}
\end{definition}

\noindent
Un $\pi$-sistem este o familie de mulțimi închisă la intersecții finite.

\end{document}
